\chapter{线性映射}\label{ch:b1:linmaps}

向量空间之间值得研究的映射，是那些与结构相容的映射：
即\emph{线性映射}。其两个基本的
子空间 --- 核与像 --- 刻画单射性与
满射性；在有限维时，秩--零化度定理
将它们的维数联系为一条守恒律。投影与对称，
随后的线性型与超平面为全章收尾。

\section{定义与初步性质}

\begin{definition}\label{def:b1:linmaps:def}
设 $E, F$ 为 $K$-向量空间。映射 $u \colon E \to F$ 称
为\emph{线性映射}\index{线性映射}，当
\[
\forall x, y \in E,\ \forall \lambda \in K, \qquad
u(x + \lambda y) = u(x) + \lambda u(y).
\]
于是 $u(0) = 0$，且 $u(\sum \lambda_i x_i) = \sum \lambda_i u(x_i)$。
线性映射的集合 $\mathcal{L}(E, F)$ 本身是一个向量空间；
线性映射的复合仍为线性映射，且复合运算关于这一对映射是双线性的。
$E \to E$ 的线性映射称为\emph{自同态}；
双射线性映射称为\emph{同构}（其逆自动是线性的）；
同构的逆 $u^{-1}$、同构的复合
皆为同构。
\end{definition}

\begin{proof}[证明逆映射是线性的]
设 $u$ 为线性双射，$y, y' \in F$，$\lambda \in K$。
令 $x = u^{-1}(y)$，$x' = u^{-1}(y')$，则
\[
u\bigl(x + \lambda x'\bigr) = u(x) + \lambda u(x')
= y + \lambda y' ,
\]
并对两端作用 $u^{-1}$：$u^{-1}(y + \lambda y') = x
+ \lambda x' = u^{-1}(y) + \lambda\,u^{-1}(y')$。这里并未对
$u^{-1}$ 作任何计算：线性性只经由
$u$ 的定义性质本身传递 --- 这一模式值得记住，
因为结构常常随双射免费而来。
\end{proof}

\begin{proposition}[线性映射由基上的值确定]\label{prop:b1:linmaps:basis}
设 $(e_1, \dots, e_n)$ 为 $E$ 的一组基，$(v_1, \dots, v_n)$
为 $F$ 中任意向量。则存在唯一的线性映射 $u \colon
E \to F$ 满足 $u(e_i) = v_i$（对所有 $i$）。此外：
\[
u \text{ 单射} \iff (v_i) \text{ 线性无关};
\qquad
u \text{ 满射} \iff (v_i) \text{ 生成 } F .
\]
\end{proposition}

\begin{proof}
存在性与唯一性：每个 $x$ 都有唯一的坐标 $x = \sum
\lambda_i e_i$（见\cref{prop:b1:vspaces:coordinates}）；由线性性
可得 $u(x) = \sum \lambda_i v_i$，且此公式确实定义了一个
线性映射。

单射性：由下文\cref{prop:b1:linmaps:kernel}，$u$
为单射当且仅当其核平凡。而 $u\bigl(\sum\lambda_i
e_i\bigr) = 0$ 恰表示 $\sum\lambda_i v_i = 0$。若 $(v_i)$
线性无关，则此式迫使每个 $\lambda_i = 0$，即核只含 $0$，
从而为单射。若 $(v_i)$ 线性相关，则一个非平凡关系
$\sum\lambda_i v_i = 0$ 给出核中的非零向量 $\sum\lambda_i e_i$
（$(e_i)$ 线性无关）：不是单射。
两个条件逐项对应。

满射性：$u$ 的像即全体 $\sum\lambda_i
v_i$，亦即恰为 $\operatorname{Vect}(v_1, \dots, v_n)$，它
等于 $F$ 当且仅当该族是生成族。
\end{proof}

\begin{definition}[核与像]\label{def:b1:linmaps:kerim}
对 $u \in \mathcal{L}(E, F)$：
\[
\ker u = \{x \in E : u(x) = 0\} \subseteq E,
\qquad
\operatorname{im} u = u(E) \subseteq F ,
\]
二者都是子空间（利用子空间判别准则，
直接验证）。\index{核}\index{像}
\end{definition}

\begin{method}[核与像的实际求法]
\label{met:b1:linmaps:kerim}
\emph{核}：将 $u(x) = 0$ 写成关于
$x$ 的坐标（或系数）的方程组，求解并参数化 --- 核
随之附带得到一组基
（见\cref{met:b1:findim:computedim}）。\emph{像}：它是
$E$ 的\emph{任一}生成族之像张成的子空间 ---
通常取一组基，于是 $\operatorname{im} u =
\operatorname{Vect}\bigl(u(e_1), \dots, u(e_n)\bigr)$；然后
消去多余的像以抽取一组基。\emph{捷径}：
计算两者中较易的一个，并由此免费得到另一者的
维数，由秩--零化度定理可知
（见\cref{thm:b1:linmaps:ranknullity}）；当已知像的某个合理候选
子空间时，比较维数可将
容易的包含升级为相等
（见\cref{thm:b1:findim:subspaces}）。两种捷径均在下文
\cref{ex:b1:linmaps:delta}中使用。
\end{method}

\begin{proposition}\label{prop:b1:linmaps:kernel}
$u$ 为单射 $\iff$ $\ker u = \{0\}$；$u$ 为满射 $\iff$
$\operatorname{im} u = F$。
\end{proposition}

\begin{proof}
与群的情形相同（见\cref{prop:b1:structures:kernel}）：$u(x) = u(y) \iff
u(x - y) = 0 \iff x - y \in \ker u$。第二个论断
即为定义。
\end{proof}

\section{秩--零化度定理}

\begin{definition}\label{def:b1:linmaps:rank}
\emph{秩}\index{秩!线性映射的}是指：对 $u \in \mathcal{L}(E,
F)$（$E$ 为有限维），其秩为 $\operatorname{rk} u = \dim
\operatorname{im} u$ --- 也是族
$\bigl(u(e_1), \dots, u(e_n)\bigr)$ 的秩，其中 $(e_i)$ 为 $E$ 的任一基。
\end{definition}

\begin{theorem}[秩--零化度定理]\label{thm:b1:linmaps:ranknullity}
设 $E$ 为有限维，$u \in \mathcal{L}(E, F)$。
则\index{秩--零化度定理}
\[
\dim E = \dim \ker u + \operatorname{rk} u .
\]
更确切地说，若 $S$ 是 $\ker u$ 在
$E$ 中的任一补子空间，则 $u$ 的限制为一个\emph{同构}，从 $S$ 映满
$\operatorname{im} u$。
\end{theorem}

\begin{proof}
设 $S$ 满足 $E = \ker u \oplus S$
（见\cref{thm:b1:findim:subspaces}），并令 $v \colon S \to
\operatorname{im} u$ 为 $u$ 的限制。

$v$ 为单射：$\ker v = S \cap \ker u = \{0\}$。

$v$ 为满射：任一 $u(x)$（$x = k + s$，$k \in \ker u$，$s
\in S$）等于 $u(s) = v(s)$。

因此 $v$ 是同构；同构将基映为基
（见\cref{prop:b1:linmaps:basis}），故 $\dim S = \dim\operatorname{im}
u$，再由 $\dim E = \dim\ker u + \dim S$ 即得结论。
\end{proof}

\begin{example}[按指定要求构造映射]
\label{ex:b1:linmaps:tospec}
构造 $u \in \mathcal{L}(\R^3)$，使 $\ker u =
\operatorname{Vect}(1,1,1)$，且 $\operatorname{im} u = \{z =
0\}$。先作检验：由秩--零化度定理需 $1 + 2 = 3$ ---
相容，故解可能存在。取一组与核相适应的基，
例如 $\bigl((1,1,1),\ e_1,\ e_2\bigr)$
（见\cref{ex:b1:findim:completion}），并规定各基向量的像
（见\cref{prop:b1:linmaps:basis}）：
\[
u(1,1,1) = 0, \qquad u(e_1) = e_1, \qquad u(e_2) = e_2 .
\]
于是 $\ker u \supseteq \operatorname{Vect}(1,1,1)$，且
$\operatorname{im} u = \operatorname{Vect}(e_1, e_2) = \{z =
0\}$；由秩--零化度定理，$\dim\ker u = 1$，故核恰为
所规定的直线。具体地，分解 $(x, y, z)
= z(1,1,1) + (x - z)e_1 + (y - z)e_2$：
\[
u(x, y, z) = (x - z,\ y - z,\ 0).
\]
此方法可推广：指定核为 $N$、像为 $I$ 的线性映射
存在，当且仅当 $\dim N + \dim I = \dim E$
--- 必要性即秩--零化度定理，充分性即本例的
构造。
\end{example}

\begin{corollary}\label{cor:b1:linmaps:samedim}
若 $\dim E = \dim F$（有限维），则对 $u \in \mathcal{L}(E, F)$：
\[
u \text{ 单射} \iff u \text{ 满射} \iff u \text{
双射}.
\]
特别地，有限维空间上的自同态满足此结论。（在无限维时此结论不成立：
在 $K[X]$ 上，求导映射是满射，
但非单射，而 $P \mapsto XP$ 是单射，
但非满射。）
\end{corollary}

\begin{proof}
单射 $\iff \dim\ker u = 0 \iff \operatorname{rk} u = \dim E =
\dim F \iff \operatorname{im} u = F$（满维子空间
即全空间，见\cref{thm:b1:findim:subspaces}）$\iff$ 满射。
\end{proof}

\begin{example}[从结构看插值]\label{ex:b1:linmaps:interpolation}
任取两两不同的 $x_0, \dots, x_n$，令 $u \colon \R_n[X] \to
\R^{n+1}$，$P \mapsto (P(x_0), \dots, P(x_n))$，显然线性。其核
为 $\{P : \deg P \leq n,\ n+1 \text{ 个根}\} = \{0\}$
（见\cref{cor:b1:poly:nroots}）。维数相等，都是 $n + 1$，故 $u$ 是
同构 --- 拉格朗日插值多项式的存在性\emph{与}唯一性
（见\cref{thm:b1:poly:lagrange}），一行即得。

同样的一行式论证也适用于混合函数值与
导数值的数据：$v \colon \R_3[X] \to \R^4$，$P \mapsto \bigl(P(0),
P'(0), P(1), P'(1)\bigr)$ 显然线性，其核由
次数 $\leq 3$ 且在 $0$ \emph{与}
$1$ 处有二重根的多项式组成，即被 $X^2(X-1)^2$（次数为 $4$）整除：只有 $P =
0$。维数再次相等：任一组数据 $(P(0),
P'(0), P(1), P'(1))$ 恰由唯一一个三次多项式实现 ---
即 埃尔米特插值，由核的计算在写出任何公式之前即已保证，
（\cref{ch:b1:det} 的周末问题
在那里遇到它的行列式）。
\end{example}

\begin{example}[秩--零化度定理的应用：差分算子]
\label{ex:b1:linmaps:delta}
设 $\Delta \colon \R_n[X] \to \R_n[X]$，$P \mapsto P(X+1) -
P(X)$，显然线性。核：若 $\Delta P = 0$，则 $P(0) = P(1) =
P(2) = \dots$，故 $P - P(0)$ 有无穷多个根，
恒为零（见\cref{cor:b1:poly:nroots}）：$\ker\Delta$ 为
常值多项式组成的直线。由秩--零化度定理，$\operatorname{rk}\Delta = (n + 1) - 1
= n$。由于 $\deg \Delta P < \deg P$ 对非常值的 $P$ 成立（最高次项
相消），$\operatorname{im}\Delta \subseteq \R_{n-1}[X]$，
其维数恰为 $n$，因此包含即为相等。
结论是，无需计算原像：\emph{任一}
多项式 $Q$（次数 $\leq n - 1$）都是差分 $Q = P(X+1)
- P(X)$ --- 离散原函数存在。（参看
\cref{ch:b1:vspaces} 的周末问题，其中 $\Delta$ 在二项式基下
被显式求逆。）
\end{example}

\begin{example}[复合映射中的秩的核算]
\label{ex:b1:linmaps:rankbookkeeping}
在 $\R_2[X]$ 上，取求导算子 $D(P) = P'$（秩为 $2$：像为 $\R_1[X]$，核为常数多项式），
与自身复合。则 $D
\circ D = D^2$ 把 $P \mapsto P''$，其像为 $\R_0[X]$：秩为
$1$。与一般的界作比较，粗略的界给出
$\operatorname{rk} D^2 \leq \min(2, 2) = 2$；而
\cref{exo:b1:linmaps:12} 中的精确公式准确地刻画了秩的损失，
\[
\operatorname{rk} D^2 = \operatorname{rk} D -
\dim\bigl(\ker D \cap \operatorname{im} D\bigr) = 2 - 1 = 1 ,
\]
由于常数多项式（外层 $D$ 的核）
包含于 $\R_1[X]$（内层 $D$ 的像）中，其维数为 $1$。秩
恰好在外层的核“撞上”内层的像之处丢失——
当复合映射的秩出现异常时，应记住这一点。
\end{example}

\begin{example}[欧拉算子的核与像]
\label{ex:b1:linmaps:euler}
在 $\R_n[X]$ 上，令 $u(P) = X\,P'$（它是线性的：求导与乘以 $X$ 皆为线性运算）。
\emph{核}：由 $XP' = 0$ 推出 $P' =
0$（多项式之积为零，当且仅当其中一个因式为零），
故 $\ker u$ 为常数多项式构成的直线。\emph{像}：
在单项式基上，
\[
u(X^k) = k\,X^{k} \qquad (k = 0, 1, \dots, n),
\]
于是 $\operatorname{im} u = \operatorname{Vect}(X, 2X^2, \dots,
nX^n) = \operatorname{Vect}(X, X^2, \dots, X^n)$：
常数项为零的多项式。与秩–零化度定理对照：
$\operatorname{rk} u = (n + 1) - 1 = n$，正是所求的
维数。值得记住的有两点。其一，这里
$\operatorname{im} u \oplus \ker u = \R_n[X]$——但这只是该算子的一个
\emph{偶然现象}，并非定理：对于类似平移的
$v(P) = P'$ 在 $\R_1[X]$ 上，$\ker v = \operatorname{im}
v = \R_0[X]$，而其和不是直和。其二，关系式
$u(X^k) = kX^k$ 说明每个单项式仅被 $u$ 作倍数的伸缩——
这是与该映射相适应的一组基，即
在本科二年级分册中展开的特征值思想的雏形。
\end{example}

\section{投影与对称}

\begin{definition}\label{def:b1:linmaps:projection}
设 $E = F \oplus G$。称 $F$ 上沿 $G$ 的
\emph{投影}\index{投影}，将 $x = f + g$（分解唯一）映为 $p(x) = f$；
与之相应的\emph{对称}\index{对称}为 $s(x)
= f - g$。两者都是线性的，且 $s = 2p - \mathrm{id}$。
\end{definition}

\begin{theorem}[代数刻画]\label{thm:b1:linmaps:projchar}
\begin{enumerate}
  \item 自同态 $p$ 是投影（即对某个 $F$ 上
        沿某个 $G$ 的投影）当且仅当 $p \circ p = p$；此时 $F =
        \operatorname{im} p = \ker(p - \mathrm{id})$，且 $G = \ker
        p$.
  \item 自同态 $s$ 为对称，当且仅当 $s \circ s
        = \mathrm{id}$；此时 $E = \ker(s - \mathrm{id}) \oplus
        \ker(s + \mathrm{id})$.
\end{enumerate}
\end{theorem}

\begin{proof}
(1) 投影满足 $p(f + g) = f$ 及 $p(f) = f$：$p^2 = p$。
反之，设 $p^2 = p$；令 $F = \operatorname{im} p$，$G = \ker
p$。每个 $x$ 均可写作 $x = p(x) + (x - p(x))$，其中 $p(x) \in F$，且
$p\bigl(x - p(x)\bigr) = p(x) - p^2(x) = 0$，因此 $E = F + G$。若 $y \in
F \cap G$：$y = p(z)$ 且 $p(y) = 0$，从而 $y = p(z) = p^2(z) = p(y) =
0$：直和，且 $p$ 为 $F$ 上沿 $G$ 的投影。
最后，在 $F$ 上：由 $y = p(z)$ 得 $p(y) = y$，故 $F \subseteq
\ker(p - \mathrm{id})$；反之，由 $p(y) = y$ 知 $y$ 落在
像中。

(2) 对应关系 $s = 2p - \mathrm{id}$，$p = \frac{s +
\mathrm{id}}2$ 是自同态之间的双射，且在其下
\[
s^2 = 4p^2 - 4p + \mathrm{id} = \mathrm{id}
\iff 4p^2 = 4p \iff p^2 = p :
\]
对称恰与投影一一对应。为考察子空间，
$s(x) = x \iff p(x) = x$，故 $\ker(s - \mathrm{id}) =
\operatorname{im} p = F$；又 $s(x) = -x \iff 2p(x) = 0 \iff x
\in \ker p = G$，故 $\ker(s + \mathrm{id}) = G$。
（1）中的直和 $E = F \oplus G$ 化为所预告的分解，
即分解为 $s$ 的不动向量与反向向量。
\end{proof}

\begin{example}[投影及其对称的显式计算]
\label{ex:b1:linmaps:projexplicit}
在 $\R^2$ 中，考虑 $F = \operatorname{Vect}(1,1)$ 上沿 $G =
\operatorname{Vect}(0,1)$ 的投影。将 $(x, y) = a(1,1) + b(0,1)$ 作分解：
第一个坐标给出 $a = x$，第二个给出 $b = y - x$。
于是
\[
p(x, y) = (x, x),
\qquad
s(x, y) = 2p(x,y) - (x,y) = (x,\ 2x - y).
\]
验算：$p(p(x,y)) = p(x,x) = (x,x)$，且 $s(s(x,y)) =
s(x, 2x - y) = (x, 2x - (2x - y)) = (x, y)$。从几何上看，$s$
是沿竖直方向、关于直线 $y = x$ 的“斜反射”，
它逐点固定 $F$，并反转 $G$。若
改为在同一 $F$ 上沿 $G' =
\operatorname{Vect}(1,-1)$ 投影，则公式变为
$p'(x,y) = \bigl(\frac{x+y}2, \frac{x+y}2\bigr)$：投影
由其像\emph{与}其核共同决定，
而绝不能仅凭像来决定。
\end{example}

\begin{omfigure}
\begin{tikzpicture}[scale=1.15]
  \draw[->, thin] (-0.4,0) -- (3.4,0) node[below] {$x$};
  \draw[->, thin] (0,-0.4) -- (0,3.8) node[left] {$y$};
  \draw[omDef, very thick] (-0.3,-0.3) -- (3.1,3.1)
    node[right, pos=0.98, font=\small] {$F:\ y = x$};
  \draw[omProp, thick, ->] (2.6,0.2) -- (2.6,0.9)
    node[right, font=\small] {$G$};
  \fill (2,0.5) circle (0.05) node[right, font=\small] {$M$};
  \fill (2,2) circle (0.05) node[left, font=\small] {$p(M)$};
  \fill (2,3.5) circle (0.05) node[left, font=\small] {$s(M)$};
  \draw[omThm, dashed] (2,0.5) -- (2,3.5);
\end{tikzpicture}

{\small 投影到 $F = \operatorname{Vect}(1,1)$ 上，沿
$G = \operatorname{Vect}(0,1)$ 的投影及其对称，取点 $M
= (2,\ 0.5)$：$M$ 竖直滑动，与 $F$ 相交于 $p(M) = (2,2)$
并落到 $s(M) = 2p(M) - M = (2,\ 3.5)$，位于 $F$ 上方，
其距离（沿 $G$ 度量）与 $M$ 的下方距离相等。}
\end{omfigure}

\section{线性型与超平面}

\begin{definition}\label{def:b1:linmaps:forms}
称\emph{线性型}\index{线性型}为 $E$ 上的线性映射
$\varphi \colon E \to K$。\emph{超平面}\index{超平面}
是 $E$（$\dim E = n$）的 $n - 1$ 维子空间。
\end{definition}

\begin{example}[求值型及其超平面]
\label{ex:b1:linmaps:evaluationform}
在 $\R_2[X]$ 上，求值 $\varphi(P) = P(2)$ 是一个线性
型，非零（$\varphi(1) = 1$）。它的核是
在 $2$ 处取零值的多项式构成的超平面，即（因式定理，
\cref{thm:b1:poly:factor}）$X - 2$ 的倍式，它们在
$\R_2[X]$ 中：
\[
\ker\varphi = \operatorname{Vect}\bigl(X - 2,\ X(X - 2)\bigr),
\qquad \dim = 2 .
\]
在基 $(1, X, X^2)$ 下取坐标，$\varphi(a + bX + cX^2) = a +
2b + 4c$。有限维空间上的每个线性型，取定基之后，
都是坐标的一个固定线性表达式——
线性型即“行向量”，正如
在\cref{ch:b1:matrices}中将明确的那样，此处的系数行
$(1, 2, 4)$ 是一个范德蒙德行：求值型
恰是\cref{ch:b1:det}的周末
问题中的插值理论进入线性代数的途径。
\end{example}

\begin{theorem}\label{thm:b1:linmaps:hyperplanes}
$E$ 的超平面恰为非零线性
型的核；两个非零线性型有相同的核，当且仅当它们
成比例。
\end{theorem}

\begin{proof}
若 $\varphi \neq 0$，则 $\operatorname{rk}\varphi = 1$（像为
$K$ 的非零子空间），故 $\dim\ker\varphi = n - 1$：是一个超平面。
反之，设 $H$ 为超平面，$(e_1, \dots, e_{n-1})$ 为
$H$ 的一组基，由 $e_n$ 补全：“最后一个坐标”这一线性型的核恰为
$H$。

成比例的线性型有相同的核。反之，设
$\ker\varphi = \ker\psi = H$，取 $a \notin H$：每个 $x$ 可写为
$x = h + \lambda a$（因 $E = H \oplus Ka$），且
\[
\varphi(x) = \lambda \varphi(a), \qquad \psi(x) = \lambda\psi(a):
\]
于是 $\varphi = \frac{\varphi(a)}{\psi(a)}\,\psi$。
\end{proof}

\begin{example}\label{ex:b1:linmaps:hyperplane}
在 $K^n$ 中，超平面是解集 $\{a_1 x_1 + \dots + a_n
x_n = 0\}$，其中 $a_i$ 不全为零，正是
$\R^3$ 中过原点的平面所熟悉的方程。在函数空间中，求值型
$P \mapsto P(1)$ 或 $f \mapsto \int_0^1 f$ 定义了
$\R_n[X]$ 与 $C(\intcc{0}{1})$ 的超平面（参见\cref{exo:b1:findim:6}）。
\end{example}

\begin{example}[同一超平面的三种处理]
\label{ex:b1:linmaps:hyperplanerun}
取 $\varphi(x, y, z) = x - 2y + 3z$（定义在 $\R^3$ 上），并设 $H =
\ker\varphi$。\emph{基}：解出 $x = 2y - 3z$：
\[
(2y - 3z,\ y,\ z) = y\,(2, 1, 0) + z\,(-3, 0, 1),
\]
两个自由向量：$\dim H = 2$，即超平面，与
\cref{thm:b1:linmaps:hyperplanes}的结论一致（由 $\varphi \neq
0$ 得出）。\emph{补直线}：$H$ 之外的任一向量张成一条补直线，
例如 $a = (1, 0, 0)$（$\varphi(a) = 1 \neq 0$）；任意向量 $v$ 的
分解是显式的：
\[
v = \underbrace{\bigl(v - \varphi(v)\,a\bigr)}_{\in\,H}
+ \underbrace{\varphi(v)\,a}_{\in\,\operatorname{Vect}(a)},
\]
因为 $\varphi\bigl(v - \varphi(v)a\bigr) = \varphi(v) -
\varphi(v)\varphi(a) = 0$。\emph{成比例}：若 $\psi(x,y,z)
= -2x + 4y - 6z$，则 $\psi = -2\varphi$，二者的核都是
$H$；反之，在 $H$ 上为零的线性型都是
$\varphi$ 的倍数（见\cref{exo:b1:linmaps:8}）——超平面
的方程在相差一个比例因子的意义下唯一，这一事实在几何中
用于研究平面时被频繁使用。
\end{example}

\begin{remark}[常见误区]
\label{rem:b1:linmaps:pitfalls}
\emph{核与像位于不同的空间中}：$\ker u
\subseteq E$，$\operatorname{im} u \subseteq F$；和 $\ker u
+ \operatorname{im} u$ 只在自同态情形下才有意义，且即便如此，
它也未必是直和（$u(x, y) = (y, 0)$ 满足 $\ker u
= \operatorname{im} u$；\cref{exo:b1:linmaps:7}给出了
直和成立的判别条件）。
\emph{$u^2 = 0$ 并不意味着 $u = 0$}：同一个 $u(x,y) =
(y, 0)$ 的平方为零而其本身并不为零——$u^2 = 0$ 真正
说明的是 $\operatorname{im} u \subseteq \ker u$
（见\cref{exo:b1:linmaps:5}）。
\emph{单射 $\iff$ 满射 需要有限且相同的
维数}：在 $K[X]$ 上，求导映射是满射而非
单射，$P \mapsto XP$ 为单射而非满射
（见\cref{cor:b1:linmaps:samedim}）；在
\emph{维数不同}的空间之间，其中一个蕴含根本不可能成立
（$\operatorname{rk} u \leq \min(\dim E, \dim F)$）。
\emph{指定像只在基上可行，对任意向量族则不行}：
要求 $u(1, 0) = a$、$u(0, 1) = b$、$u(1, 1) = c$
则 $u$ 被过度确定，除非 $c = a + b$；线性映射在基上可以任意指定，
在其他处则受约束。
\emph{秩在复合下不保持不变}：它只会下降，
$\operatorname{rk}(vu) \leq \min(\operatorname{rk} u,
\operatorname{rk} v)$（见\cref{exo:b1:linmaps:4}），其精确损失
在\cref{exo:b1:linmaps:12}中度量。
\end{remark}

\begin{remark}[这些映射的用途]
\label{rem:b1:linmaps:whereused}
线性映射即将化为\emph{矩阵}：一旦基
确定，在\cref{ch:b1:matrices}中，每个 $u \in \mathcal{L}(E,
F)$ 都由一个矩形数组刻画，复合化为矩阵
乘积 --- 秩–零度定理随后推动线性方程组
理论，见\cref{ch:b1:det}。投影将在
\cref{ch:b1:euclid}中以其最有用的特例再次出现，即
\emph{正交}投影，其中核取为
与像垂直的子空间。下文的周末问题将
投影算子的代数推到一年级工具所能及的极限，直至
菲廷引理；本科二年级分册将借助迹走得更远，
并结合特征值理论，其中投影到稳定子空间的投影算子
是基本构件。
\end{remark}

\begin{remark}[第 3 册中的延伸视角：秩–零度定理再现三次]
\label{rem:b1:linmaps:perspectives}
守恒律 $\dim E = \dim\ker u + \operatorname{rk} u$
将在本册结束前被重新审视三次。在
\cref{ch:b1:det}中，它化为解集的形状：一个
相容方程组含 $p$ 个未知量、秩为 $r$，其解集
维数为 $p - r$ --- 即伪装起来的核维数。在
\cref{ch:b1:euclid}中，它正交分解为 $\dim F + \dim
F^\perp = \dim E$，并支撑每一次距离计算。在
\cref{ch:b1:multivar}的周末问题中，它是
最小二乘的“记账员”：$n$ 个观测、$2$ 个拟合
参数、$n - 2$ 个残差维数，而勾股
恒等式 $\norm b^2 = \norm p^2 + \norm{b - p}^2$ 是
秩–零度定理在欧几里得意义下的影子。一个定理，四副装束。
\end{remark}

\section{练习}

\begin{exercise}[$\star$]\label{exo:b1:linmaps:1}
哪些映射是线性映射？
\begin{enumerate}
  \item $\R^2 \to \R^2$, $(x, y) \mapsto (x + y, x - 2y)$;
  \item $\R^2 \to \R$, $(x, y) \mapsto xy$;
  \item $\R[X] \to \R[X]$, $P \mapsto P' + XP$;
  \item $\mathcal{F}(\R,\R) \to \R$, $f \mapsto f(3)$.
\end{enumerate}
\end{exercise}

\begin{exercise}[$\star$]\label{exo:b1:linmaps:2}
设 $u \colon \R^3 \to \R^3$，$(x,y,z) \mapsto (x + y - z,\; 2x + y
+ z,\; 3x + 2y)$。求 $\ker u$（基、维数），
$\operatorname{rk} u$，以及 $\operatorname{im} u$ 的一组基。$u$ 是否
为单射？是否为满射？
\end{exercise}

\begin{exercise}[$\star$]\label{exo:b1:linmaps:3}
设 $u \colon \R_n[X] \to \R_n[X]$，$P \mapsto P - P'$。证明
$u$ 是同构：一种方法借助 $\ker u$，另一种方法通过给出
逆映射\emph{（考虑 $P + P' + P'' + \dots$）}。
\end{exercise}

\begin{exercise}[$\star$]\label{exo:b1:linmaps:4}
设 $u \in \mathcal{L}(E, F)$ 与 $v \in \mathcal{L}(F, G)$，各空间
均为有限维。证明：
\[
\operatorname{rk}(v \circ u) \leq
\min\bigl(\operatorname{rk} u,\ \operatorname{rk} v\bigr).
\]
\end{exercise}

\begin{exercise}[$\star\star$]\label{exo:b1:linmaps:5}
设 $u$ 为 $E$ 上的自同态（有限维），满足 $u^2 =
0$。证明 $\operatorname{im} u \subseteq \ker u$，从而
$\operatorname{rk} u \leq \frac{\dim E}{2}$。对 $E = \R^2$，给出
一个取等号的例子。
\end{exercise}

\begin{exercise}[$\star\star$]\label{exo:b1:linmaps:6}
设 $p, q$ 为 $E$ 上的投影，满足 $p \circ q = q \circ p$。
证明 $p \circ q$ 是投影，且
\[
\operatorname{im}(pq) = \operatorname{im} p \cap \operatorname{im}
q ,
\qquad
\ker (pq) = \ker p + \ker q .
\]
\end{exercise}

\begin{exercise}[$\star\star$]\label{exo:b1:linmaps:7}
设 $u \in \mathcal{L}(E)$，$E$ 为有限维。证明以下
三个条件等价：
\begin{enumerate}
  \item $E = \ker u \oplus \operatorname{im} u$;
  \item $\ker u = \ker u^2$;
  \item $\operatorname{im} u = \operatorname{im} u^2$.
\end{enumerate}
\end{exercise}

\begin{exercise}[$\star\star$]\label{exo:b1:linmaps:8}
设 $\varphi, \psi$ 为 $E$ 上的线性型，满足 $\ker\varphi
\subseteq \ker\psi$。证明 $\psi = \lambda\varphi$，其中
$\lambda \in K$ 为某个数（包括退化情形）。
\end{exercise}

\begin{exercise}[$\star\star\star$]\label{exo:b1:linmaps:9}
设 $u \in \mathcal{L}(E)$，$\dim E = n$，且 $u^n = 0$
但 $u^{n-1} \neq 0$（即\emph{极大幂零}自同态）。
取 $x$ 使 $u^{n-1}(x) \neq 0$；证明 $\bigl(x, u(x), \dots,
u^{n-1}(x)\bigr)$ 是 $E$ 的一组基。\emph{（将 $u$ 的各次幂作用于一个
零线性组合上，从 $u^{n-1}$ 开始。）}
\end{exercise}

\begin{exercise}[$\star\star\star$]\label{exo:b1:linmaps:10}
设 $f \in \mathcal{L}(\R^n)$，满足 $f \circ f = -\mathrm{id}$。
\begin{enumerate}
  \item 证明 $f$ 是同构，且不存在 $x \neq 0$
        满足 $f(x) = \lambda x$，其中 $\lambda \in \R$。
  \item 证明 $n$ 为偶数。\emph{提示：取 $x_1 \neq 0$；证明
        $\operatorname{Vect}(x_1, f(x_1))$ 是一个在
        $f$ 下稳定的平面；取 $x_2$ 在其外并迭代，证明
        $\bigl(x_1, f(x_1), x_2, f(x_2), \dots\bigr)$ 仍线性无关。}
\end{enumerate}
\end{exercise}

\begin{exercise}[$\star\star$]\label{exo:b1:linmaps:11}
设 $u, v \in \mathcal{L}(E, F)$，且空间均为有限维。
证明下面的双侧估计：
\[
\abs{\operatorname{rk} u - \operatorname{rk} v}
\;\leq\; \operatorname{rk}(u + v)
\;\leq\; \operatorname{rk} u + \operatorname{rk} v .
\]
\emph{（对上界，比较 $\operatorname{im}(u+v)$ 与
$\operatorname{im} u + \operatorname{im} v$；对下界，
巧妙地利用上界。）}
\end{exercise}

\begin{exercise}[$\star\star\star$]\label{exo:b1:linmaps:12}
（弗罗贝尼乌斯不等式）设 $u \in \mathcal{L}(E, F)$，$w \in
\mathcal{L}(F, G)$ 与 $v \in \mathcal{L}(G, H)$，各空间均为
有限维。证明如下精确公式
\[
\operatorname{rk}(v \circ w) = \operatorname{rk} w -
\dim\bigl(\ker v \cap \operatorname{im} w\bigr),
\]
并由此推出弗罗贝尼乌斯不等式
\[
\operatorname{rk}(v \circ w) + \operatorname{rk}(w \circ u)
\;\leq\; \operatorname{rk} w + \operatorname{rk}(v \circ w \circ
u) .
\]
验证 $w = \mathrm{id}_F$ 时的情形即为西尔维斯特
不等式，其矩阵形式的证明见\cref{exo:b1:matrices:10}。
\end{exercise}

\section{问题：投影的运算与费廷引理}

\begin{problem}\label{pb:b1:linmaps:1}
投影即直和所产生的自同态，反之亦然：
每个分解 $E = F_1 \oplus \dots \oplus F_k$ 本质上都是
一族和为恒等映射的投影。本题展开这一对应——
即单个投影的代数，
两个投影、$k$ 个投影的代数，
再将同样的稳定化思想用于任意自同态，并证明
\emph{费廷引理}：每个有限维空间上的
自同态都可分解为一个幂零部分与一个
可逆部分。全文中 $E$ 为 $K$-向量空间，其
维数为 $n$，而\emph{投影}指 $p \in \mathcal{L}(E)$ 且
满足 $p^2 = p$（见\cref{thm:b1:linmaps:projchar}）。
\index{投影}\index{费廷引理}\index{幂零}

\medskip
\noindent\textbf{第一部分 --- 单个投影的代数。}
设 $p$ 为投影，$p \neq 0$，$p \neq \mathrm{id}$。
\begin{enumerate}
  \item 证明 $\mathrm{id} - p$ 是投影，并求出
        $\operatorname{im}(\mathrm{id} - p)$ 与
        $\ker(\mathrm{id} - p)$。
  \item 计算 $(\lambda\,\mathrm{id} + \mu\,p)^2$，并确定
        全部数对 $(\lambda, \mu) \in K^2$，
        使 $\lambda\,\mathrm{id} + \mu\,p$ 为投影。
  \item 证明平面 $\operatorname{Vect}(\mathrm{id}, p)$
        （属于 $\mathcal{L}(E)$）在复合下稳定，并证明
        对任一多项式 $Q \in K[X]$，有
        \[
        Q(p) = Q(0)\,\mathrm{id} + \bigl(Q(1) -
        Q(0)\bigr)\,p .
        \]
  \item 确定对哪些 $(\lambda, \mu)$ 而言，映射
        $\lambda\,\mathrm{id} + \mu\,p$ 可逆，并给出
        其逆写成 $\alpha\,\mathrm{id} + \beta\,p$ 的形式。
        并通过
        $\lambda\,\mathrm{id} + \mu\,p$ 在 $\operatorname{im} p$
        与 $\ker p$ 上的作用来解释结果。
  \item 设 $p'$ 为另一投影，且具有\emph{相同的像}
        即 $\operatorname{im} p' = \operatorname{im} p$。证明
        $p\,p' = p'$ 与 $p'\,p = p$。这些等式
        对同一子空间、沿不同核所作投影的复合，
        说明了什么？
\end{enumerate}

\medskip
\noindent\textbf{第二部分 --- 两个投影。}设 $p, q$ 为
$E$ 的投影，并假设特征不为 $2$（对
$K = \R, \C$ 成立）。
\begin{enumerate}[resume]
  \item 设 $p + q$ 为投影。展开 $(p + q)^2$，
        证明 $pq + qp = 0$；分别左乘、右乘 $p$，
        由此推出 $pq = qp$，并得出 $pq = qp =
        0$。
  \item 反之，设 $pq = qp = 0$，证明 $p + q$ 为
        投影，且
        \[
        \operatorname{im}(p + q) = \operatorname{im} p \oplus
        \operatorname{im} q,
        \qquad
        \ker(p + q) = \ker p \cap \ker q .
        \]
  \item 证明：$p - q$ 为投影当且仅当 $pq = qp
        = q$。\emph{（对 $\mathrm{id} - p$
        与 $q$ 应用第 6--7 题。）}
  \item 证明 $pq = qp = q$ 的几何意义：该等式成立
        当且仅当 $\operatorname{im} q \subseteq
        \operatorname{im} p$ 且 $\ker p \subseteq \ker q$。（此时记为 $q \leq p$，
        “$q$ 投影到较小的像上，沿
        较大的核”。）
  \item 现设 $p, q$ 可交换。回顾
        \cref{exo:b1:linmaps:6}：$pq$ 是到
        $\operatorname{im} p \cap \operatorname{im} q$ 上、沿
        $\ker p + \ker q$ 的投影。证明 $r = p + q - pq$ 为
        投影，且
        \[
        \operatorname{im} r = \operatorname{im} p +
        \operatorname{im} q,
        \qquad
        \ker r = \ker p \cap \ker q .
        \]
        \emph{（考虑 $\mathrm{id} - r = (\mathrm{id} -
        p)(\mathrm{id} - q)$。）}
\end{enumerate}

\medskip
\noindent\textbf{第三部分 --- 恒等映射的分解。}
\begin{enumerate}[resume]
  \item 设 $E = F_1 \oplus \dots \oplus F_k$，对 $x = x_1 +
        \dots + x_k$（唯一分解，$x_i \in F_i$），令
        $p_i(x) = x_i$。证明每个 $p_i$ 都是投影，且
        $p_i p_j = 0$（$i \neq j$），并且 $p_1 + \dots + p_k
        = \mathrm{id}$；确定 $\operatorname{im} p_i$ 与
        $\ker p_i$。
  \item 反之，设 $p_1, \dots, p_k \in \mathcal{L}(E)$
        满足 $p_1 + \dots + p_k = \mathrm{id}$，且 $p_i p_j =
        0$，对一切 $i \neq j$。证明每个 $p_i$ 都是
        投影，并且 $E = \operatorname{im} p_1 \oplus
        \dots \oplus \operatorname{im} p_k$。
  \item 对满足 $p + q = \mathrm{id}$ 的两个投影，证明 $pq
        = qp = 0$ 自动成立。
  \item 对满足 $p + q + r = \mathrm{id}$ 的三个投影，证明
        $p + q$ 是投影，并由第 6 题推出
        \emph{所有}两两乘积全为零 --- 从而 $E =
        \operatorname{im} p \oplus \operatorname{im} q \oplus
        \operatorname{im} r$，无需对乘积作任何假设。
  \item 对 $k$ 个投影，满足 $p_1 + \dots + p_k =
        \mathrm{id}$，先证对任意子空间，$\dim
        (F_1 + \dots + F_k) \leq \dim F_1 + \dots + \dim F_k$，
        等号成立当且仅当该和为直和；然后
        证明 $E = \operatorname{im} p_1 + \dots +
        \operatorname{im} p_k$，并证明：\emph{若}
        此外 $\sum_i \operatorname{rk} p_i \leq n$，则和
        为直和，且 $p_i p_j = 0$（$i \neq j$）。
\end{enumerate}

\medskip
\noindent\textbf{第四部分 --- 迭代核：费廷引理。}
设 $u \in \mathcal{L}(E)$，$\dim E = n$。
\begin{enumerate}[resume]
  \item 证明下面两条链对每个 $k \geq 0$ 都成立：
        \[
        \ker u^k \subseteq \ker u^{k+1},
        \qquad
        \operatorname{im} u^{k+1} \subseteq \operatorname{im}
        u^k .
        \]
  \item 证明：若 $\ker u^{r} = \ker u^{r+1}$ 对某个 $r$ 成立，
        则 $\ker u^{k} = \ker u^{r}$ 对所有 $k \geq r$ 成立；陈述
        并证明像上相应的稳定化结论。
  \item 由此推出：存在最小的整数 $r$，使得 $\ker
        u^{r} = \ker u^{r+1}$，$r \leq n$，且像
        在同一 $r$ 处稳定。
  \item （费廷引理）证明：
        \[
        E \;=\; \ker u^{r} \,\oplus\, \operatorname{im} u^{r} .
        \]
  \item 证明这两个子空间都是 $u$ 的不变子空间，
        $u$ 在 $\ker u^{r}$ 上的限制是幂零的，并且
        $u$ 在 $\operatorname{im} u^{r}$ 上的限制是
        $\operatorname{im} u^{r}$ 上的同构：于是每个
        自同态在典范直和上都是“幂零
        加可逆”的。
  \item 设 $\pi$ 为到 $\ker u^{r}$ 上、
        沿 $\operatorname{im} u^{r}$ 的投影。证明 $\pi \circ u = u
        \circ \pi$。
\end{enumerate}

\medskip
\noindent\textbf{第五部分 --- 一个算例与综合。}
\begin{enumerate}[resume]
  \item 设 $u(x, y, z) = (y, 0, z)$ 为 $\R^3$ 上的映射。计算 $u^2$ 与
        $u^3$，求稳定化指标 $r$、
        子空间 $\ker u^{r}$ 与 $\operatorname{im} u^{r}$、
        费廷投影 $\pi$，并用公式验证
        $\pi u = u\pi$，以及 $u$ 在一个因子上幂零、
        在另一个因子上双射。
  \item 证明下列等价关系：$u$ 幂零 $\iff$ $\ker u^{r} =
        E$ $\iff$ $\pi = \mathrm{id}$；并由此推出
        $n$ 维空间上的幂零自同态总满足
        $u^{n} = 0$（幂零指数永不超过
        维数）。
  \item （唯一性）设 $E = A \oplus B$，$A, B$ 均
        关于 $u$ 不变，$u|_A$ 幂零，$u|_B$
        是双射。证明 $A = \ker u^{r}$，$B =
        \operatorname{im} u^{r}$：费廷分解是
        唯一的。
  \item 综合题，用四句话作答：第三部分建立了何种
        直和与投影族
        之间的对应；为何第 14 题不需要乘积假设，
        而第 15 题需要秩的假设
        （借助本科二年级的迹可将其去除）；在何种
        意义上费廷引理是
        \cref{exo:b1:linmaps:7} 的稳定化版本；以及费廷分解的两个因子
        在本科二年级分册的特征值理论中分别变成什么；
        并说出第四部分所证明的定理。
\end{enumerate}
\end{problem}
